\documentclass[12pt, twoside]{article}
%\documentclass[12pt, twoside]{article}
\usepackage[letterpaper, margin=1in, headsep=0.2in]{geometry}
\setlength{\headheight}{0.6in}
%\usepackage[english]{babel}
\usepackage[utf8]{inputenc}
\usepackage{microtype}
\usepackage{amsmath}
\usepackage{amssymb}
%\usepackage{amsfonts}
\usepackage[nomessages]{fp} %\FPeval{\var-name}{2*sin(pi/6)}
\usepackage{siunitx} %units in math. eg 20\milli\meter
\usepackage{yhmath} % for arcs, overparenth command
\usepackage{tikz} %graphics
\usetikzlibrary{quotes, angles, arrows, arrows.meta}
\usepackage{pgfplots}
\usepgfplotslibrary{fillbetween}
\pgfplotsset{compat=1.16}
\usepackage{graphicx} %consider setting \graphicspath{{images/}}
\usepackage{parskip} %no paragraph indent
\usepackage{enumitem}
\usepackage{multicol}
\usepackage{venndiagram}

\usepackage{fancyhdr}
\pagestyle{fancy}
\fancyhf{}
\renewcommand{\headrulewidth}{0pt} % disable the underline of the header
\raggedbottom
\hfuzz=2mm %suppresses overfull box warnings

\usepackage{hyperref}
\setlength{\parskip}{0.3\baselineskip}

\fancyhead[LE]{\thepage}
\fancyhead[RO]{Markscheme}
\fancyhead[LO]{La Scuola d'Italia / Dr.\ Huson / IB Math 12 \\* 10 November 2026}

\begin{document}
\subsubsection*{1.12 Classwork: Simpson's Paradox --- Markscheme}

\emph{Lesson problem set, not scored: model answers and teaching notes.
Correlation descriptors follow the textbook scale: $\lvert r \rvert < 0.25$
very weak, $0.25$--$0.5$ weak, $0.5$--$0.75$ moderate, $0.75$--$0.95$ strong.
Data and figures: \texttt{claude-ib-project/output/correlation\_claude/}
(data set D).}

\begin{enumerate}[itemsep=0.3cm]

%--- Problem 1: Reading the two time series; currency conversion in and out of dollars; effect of a weaker and a stronger dollar on the investor ---
\item[1.]

\textbf{(a)} $100\,000 \times 1.0622 = 106\,220$ US dollars.

\textbf{(b)} $\dfrac{106\,220}{1.5759} = 67\,402.75\ldots$, so EUR $67\,400$.

percentage change $= \dfrac{67\,402.75\ldots - 100\,000}{100\,000} \times 100
= -32.6\%$ \quad $(-32.597\ldots)$

\emph{The dollar weakened by about a third against the euro. Common error:
multiplying by $1.5759$ instead of dividing (USD $\to$ EUR divides by
EUR/USD), which gives a ``gain'' to EUR $167\,000$.}

\textbf{(c)} EUR/USD fell from $1.5759$ to $1.2744$, so the dollar became
\emph{stronger}.

$\dfrac{106\,220}{1.2744} = 83\,349.02\ldots$, so EUR $83\,300$ (accept
$83\,349$).

\emph{Teaching note: in four months the same dollars gained about $24\%$ in
euro terms ($83\,349 / 67\,403$). In the autumn 2008 crisis, investors
worldwide bought dollars and US Treasuries as a safe haven.}

\textbf{(d)} The yield fell from $4.01\%$ to $2.42\%$, a fall of $1.59$
percentage points. When yields fall, the prices of existing bonds rise, so
this was good for the investor.

\emph{Accept ``fell steeply in December 2008''. Late 2008 was good for the
investor on both counts: bond prices up and the dollar stronger.}

\bigskip \hrule \bigskip

%--- Problem 2: Correlation of levels: r for each period and the whole period; Simpson's paradox explained by the shift in the period means ---
\item[2.]

\textbf{(a)} 2003--2007: $r = 0.39$, weak positive.

2008--2011: $r = 0.46$, weak positive.

2003--2011: $r = -0.30$, weak negative.

\emph{Accept ``moderate'' for $0.46$. The point to draw out: both periods
are positive, the whole is negative.}

\textbf{(b)} From the first period to the second, the mean yield fell
($4.40 \to 3.23$) and the mean EUR/USD rose ($1.249 \to 1.396$). So the
2008--2011 cloud of points lies up and to the left of the 2003--2007 cloud.
A line fitted to all the points mostly joins the two clouds, and it slopes
down. This shift between the periods outweighs the upward slope inside
each period.

\emph{Teaching background (law of total covariance): the whole-period
covariance is the weighted within-period covariance plus a between-period
term $\sum_k w_k(\bar{x}_k - \bar{x})(\bar{y}_k - \bar{y})$. Here the
between-period term is negative and dominates.}

\textbf{(c)} Within each period the relationship is the opposite: months
with higher yields had a higher EUR/USD, a \emph{weaker} dollar. The
negative whole-period value comes from the difference between the two
periods, not from how the series move together month by month. Yields and
the exchange rate both drifted over the years, so a correlation of
\emph{levels} mainly reflects those drifts and changes with the dates
chosen.

\emph{Also accept: correlation does not show causation; a third factor (the
2008 crisis and the Federal Reserve's rate cuts) drove both.}

\bigskip \hrule \bigskip

%--- Problem 3: Monthly changes: change in percentage points and percentage change; r of changes in each period; interpretation as a partial hedge for the investor ---
\item[3.]

\textbf{(a)}

\begin{center}
\setlength{\tabcolsep}{5pt}\renewcommand{\arraystretch}{1.15}
\begin{tabular}{|l|c|c|c|c|c|c|}
\hline
Month (2008) & Jul & Aug & Sep & Oct & Nov & Dec \\
\hline
change in yield (pp) & & $-0.12$ & $-0.20$ & $+0.12$ & $-0.28$ & $-1.11$ \\
\hline
change in EUR/USD (\%) & & $-5.10$ & $-4.10$ & $-7.50$ & $-3.93$ & $+6.02$ \\
\hline
\end{tabular}
\end{center}

\emph{For example, September: $\dfrac{1.4342 - 1.4955}{1.4955} \times 100
= -4.0990\ldots$. Common errors: percentage change of the yield
($-0.20/3.89 = -5.14\%$) instead of the change in percentage points;
dividing by this month's value instead of last month's.}

\textbf{(b)} 2003--2007: $r = -0.25$; 2008--2011: $r = -0.15$; 2003--2011:
$r = -0.18$. All three are weak or very weak and \emph{negative}. Unlike the
levels in 2(a), the changes give the same sign for each period and for the
whole period: there is no Simpson's paradox.

\textbf{(c)} $r < 0$: in a month when the US yield rises, EUR/USD tends to
fall, so the dollar tends to get \emph{stronger}. A stronger dollar raises
the euro value of the investor's dollars. This currency gain partly offsets
the loss on the bonds, so the two risks partly cancel, and the investor's
total risk is smaller than it would be if the two moved independently.

\emph{The offset is small: $r = -0.18$ is very weak. Class examples from
the table: October 2008 (yield $+0.12$, dollar $7.5\%$ stronger) and
December 2008 (yield $-1.11$, bond gain; dollar $6\%$ weaker, currency
loss).}

\bigskip \hrule \bigskip

%--- Problem 4: Six-point Simpson's paradox on the GDC; why changes resist the paradox; why correlating price levels misleads in a research paper ---
\item[4.]

\textbf{(a)} Period 1: $r = -1$. \quad Period 2: $r = -1$. \quad All six
points: $r = 0.807$ \quad $(0.80722\ldots)$

\emph{By hand: $\bar{x} = 4.5$, $\bar{y} = 6.5$, $S_{xy} = 33.5$,
$S_{xx} = S_{yy} = 41.5$, so $r = 33.5/41.5$.}

\textbf{(b)} Each period lies on a straight line with gradient $-1$, but
period 2 lies up and to the right of period 1: the means are $(2, 4)$ and
$(7, 9)$. Taken together, the points rise from the first group to the
second, so $r$ for all six points is positive.

\textbf{(c)} Simpson's paradox appears when the groups have different
means, so the clouds sit in different places. Monthly changes average close
to zero in every period, so the clouds of changes sit on top of each other.
There is no shift between the periods to create a trend of its own, and
the whole-period $r$ reflects the month-to-month relationship found inside
each period.

\textbf{(d)} Model answer: ``Correlate the monthly changes (returns) of the
two prices, not the prices themselves: prices trend, so the correlation of
levels mostly measures shared trends and changes sign with the dates you
choose. Report $r$ for the returns over the whole period and in each half,
to check that the relationship is stable.''

\emph{Teaching note: for gold and oil, 2015--2025, the monthly returns give
$r \approx -0.04$ for the whole period and for each half.}

\end{enumerate}
\end{document}
